\section{A quadratic root-of-unity obstruction}\label{sec:quadratic-family}

For $n\geq1$, let $R_n$ be the self-dual commutative rank-four based ring
whose fusion matrices in the basis $(1,X,Y,Z)$ are
\begingroup\small
\[
N_X=\begin{pmatrix}0&1&0&0\\1&n-1&n-1&n\\0&n-1&n&n\\0&n&n&n\end{pmatrix},\quad
N_Y=\begin{pmatrix}0&0&1&0\\0&n-1&n&n\\1&n&n&n\\0&n&n&n+1\end{pmatrix},\quad
N_Z=\begin{pmatrix}0&0&0&1\\0&n&n&n\\0&n&n&n+1\\1&n&n+1&n+1\end{pmatrix}.
\]
\endgroup
Put $\delta=(3n+\sqrt{9n^2+4})/2$. Direct substitution gives the dimension
vector $(1,\delta-1,\delta,\delta+1)$ and
\[
 \delta^2=3n\delta+1,\qquad D=3+3\delta^2=6+9n\delta.
\]

\begin{lemma}\label{lem:odd-quadratic-length}
Let $d>1$ be squarefree and let $f$ be the conductor of $\mathbb Q(\sqrt d)$:
$f=d$ if $d\equiv1\pmod4$ and $f=4d$ otherwise. Put
$\kappa=\sqrt{f/d}$ and $\lambda=(\varphi(f)-|\mu(f)|)/\kappa$.
Here $\varphi$ and $\mu$ denote Euler's totient and the M\"obius function.
If a sum of $k$ roots of unity equals $a+b\sqrt d$, with
$a,b\in\mathbb Q$, then $k\geq-a-\lambda b$.
\end{lemma}
\begin{proof}
Let $K=\mathbb Q(\sqrt d)$, let $L$ be a common cyclotomic overfield, and
write $T=\operatorname{Tr}_{L/K}/[L:K]$. Define
$\ell(a+b\sqrt d)=-a-\lambda b$. We prove $\ell(T(\xi))\leq1$
for every root of unity $\xi$.

If $m$ is the exact order of $\xi$ and $f\nmid m$, then
$K\cap\mathbb Q(\xi)=\mathbb Q$, because the conductor of $K$ is $f$.
Consequently $T(\xi)=\mu(m)/\varphi(m)$, whose absolute value is at most one.
If $m=fh$, let $\chi$ be the primitive quadratic character modulo $f$ and
write $\xi=\zeta_m^r$, with $(r,m)=1$. The induced quadratic Gauss sum gives
\[
 T(\xi)=
 \begin{cases}
 \displaystyle\frac{\mu(f)\mu(h)+\chi(r)\mu(h)\chi(h)\kappa\sqrt d}
                      {\varphi(f)\varphi(h)},&(f,h)=1,\\[6pt]
 0,&(f,h)>1.
 \end{cases}
\]
For clarity, this follows by splitting the units modulo $m$ into the two
values of $\chi$: their sum over $\chi=1$ is half the sum of the Ramanujan
sum $\mu(m)$ and the induced Gauss sum. The latter equals
$\chi(r)\mu(h)\chi(h)\sqrt f$ when $(f,h)=1$ and vanishes otherwise;
an additional power of a prime dividing $f$ cancels upon summing over its
last prime-power layer. The same formula includes nonsquarefree $h$ through
$\mu(h)=0$.
In the nonzero case, the definition of $\lambda$ gives
\[
 |\ell(T(\xi))|
 \leq\frac{|\mu(f)|+\lambda\kappa}{\varphi(f)\varphi(h)}\leq1.
\]
Applying $\ell\circ T$ to the sum proves the lemma.
\end{proof}

\subsection{The sharp conductor-forty boundary}
\begin{lemma}[A root-of-unity length certificate]\label{lem:rootlength40}
Suppose that $\xi_1,\ldots,\xi_n$ are roots of unity and
\[
 \sum_{j=1}^n\xi_j=a+b\sqrt{10},\qquad a,b\in\mathbb Q.
\]
Then $n\geq -a-8b$. If equality holds, every $\xi_j$ is either $-1$ or
is a primitive fortieth root $\zeta_{40}^{r}$ with
$\chi_{40}(r)=-1$, where $\chi_{40}(r)=(40/r)$ is the primitive quadratic
Dirichlet character of conductor $40$.
Consequently, if
\[
 n=18,\qquad\sum_{j=1}^{18}\xi_j=-2-2\sqrt{10},
\]
then exactly two of the roots are $-1$, and the normalized absolute trace of
$\sum_j\xi_j^2$ is $2$.
\end{lemma}

\begin{proof}
Put $s=\sqrt{10}$. Let $K=\mathbb Q(s)$, whose conductor is $40$, and choose a cyclotomic field
$L$ containing $K$ and all the roots under consideration. Write
\[
 T(u)=\frac{\operatorname{Tr}_{L/K}(u)}{[L:K]},\qquad
 \ell(a+bs)=-a-8b.
\]
We claim that $\ell(T(\xi))\leq1$ for every root of unity $\xi$, with precisely
the equality cases stated above. The normalized relative trace does not depend
on the choice of common cyclotomic overfield.

Let $m$ be the exact order of $\xi$. If $40\nmid m$, the intersection
$K\cap\mathbb Q(\xi)$ is $\mathbb Q$: otherwise the conductor of $K$ would
divide $m$. Restriction of automorphisms consequently gives
\[
 T(\xi)=\frac{\mu(m)}{\varphi(m)}.
\]
It follows that $\ell(T(\xi))\leq1$, with equality only when $m=2$.

Now suppose $m=40h$, and write $\xi=\zeta_m^r$, $(r,m)=1$. Since $\mu(m)=0$,
the rational coefficient of $T(\xi)$ is zero. The quadratic coefficient is
computed by the induced quadratic Gauss sum:
\[
 G_m(r)=\sum_{v\in(\mathbb Z/m\mathbb Z)^\times}
             \chi_{40}(v)\zeta_m^{rv}.
\]
For completeness, the standard induced-Gauss-sum identity gives
\[
 G_m(r)=
 \begin{cases}
 \chi_{40}(r)\mu(h)\chi_{40}(h)\,2s,&(h,40)=1,\\
 0,&(h,40)>1.
 \end{cases}
\]
Here $\mu(h)=0$ already covers nonsquarefree $h$. This identity can also be
obtained directly from the Chinese remainder theorem when $(h,40)=1$:
the sum factors into the primitive Gauss sum
$G_{40}(1)=\sqrt{40}=2s$ and the Ramanujan sum $\mu(h)$, with the displayed
quadratic character factors. If an extra power of $2$ or $5$ occurs in $m$,
summing over the last such prime-power layer gives zero.
The terms with $\chi_{40}(v)=1$ are the automorphisms fixing $K$; their sum
is $(\mu(m)+G_m(r))/2$. Since there are $\varphi(m)/2$ such automorphisms,
\[
 T(\xi)=\frac{G_m(r)}{\varphi(m)}
 =\frac{\chi_{40}(r)\mu(h)\chi_{40}(h)}{8\varphi(h)}s
 \quad\text{if }(h,40)=1,
\]
and it is zero otherwise. Thus $|\ell(T(\xi))|\leq1$. Equality with value
$1$ forces $\varphi(h)=1$ and $(h,40)=1$, hence $h=1$, and then
$\chi_{40}(r)=-1$. This proves the claim. Applying $\ell\circ T$ to the sum
of the roots proves the inequality and its equality statement.

In the particular equality case, primitive fortieth roots have normalized
absolute trace $\mu(40)/\varphi(40)=0$, whereas $-1$ has trace $-1$.
The absolute trace of $-2-2s$ is $-2$, so exactly two roots are $-1$.
The square of a primitive fortieth root is primitive of order $20$ and has
normalized absolute trace $\mu(20)/\varphi(20)=0$. The two other squares
are $1$, proving the last assertion.
\end{proof}


\begin{proposition}\label{prop:quadratic-family-n3-n4}
Let $n\geq2$, write $9n^2+4=t^2d$ with $d$ squarefree, and define
$\lambda$ as in Lemma~\ref{lem:odd-quadratic-length}.
If $n=2$, or if
\begin{equation}\label{eq:rootlength-parameter}
 n^2+nt\lambda>8n^2+4,
\end{equation}
then $R_n$ admits no spherical categorification over $\mathbb C$.
In particular this excludes every $R_n$ with $2\leq n\leq127$,
which is exactly the part of this family of multiplicity at most $128$
apart from $R_1$.
\end{proposition}
\begin{proof}
We first derive a necessary root-of-unity identity for a pseudo-unitary
categorification of $R_n$. Give such a category its canonical spherical
structure, and use $F,I$ for forgetful functor and central induction.
The four characters, displayed by rows, are
\[
 \begin{pmatrix}
 1&\delta-1&\delta&\delta+1\\
 1&\delta'-1&\delta'&\delta'+1\\
 1&0&1&-1\\
 1&1&-1&0
 \end{pmatrix},\qquad \delta'=-\delta^{-1}.
\]
The formal codegrees are $D,D'=3+3(\delta')^2,3,3$. Accordingly,
\[
 I(1)=1\oplus A\oplus B\oplus C,
 \qquad d_A=D/D'=1+3n\delta,\quad d_B=d_C=D/3=2+3n\delta.
\]
Every summand has multiplicity one and twist one. Comparing the coefficients
of $1$ and $\delta$ in their restriction dimensions shows that their columns are
\[
 (1,a,3n-2a,a),\quad
 (1,b,3n-1-2b,b+1),\quad
 (1,c,3n-1-2c,c+1).
\]
The first row of $MM^{\mathsf T}=Q:=\sum_U N_U^2$, for the full induction
matrix $M$, gives $a+b+c=3n-1$.

Here the low part is forced without searching over $a,b,c$. Put
$v=(0,1,-2,1)^{\mathsf T}$. Multiplication of the displayed matrices gives
$v^{\mathsf T}Qv=18$. On the three low columns, the values of $v$ are
\[
 6(a-n),\quad 6(b-n)+3,\quad6(c-n)+3.
\]
Their squares sum to at most $18$. The last two are odd multiples of three,
so their squares already sum to at least $18$. It follows that
$a=n$ and $\{b,c\}=\{n-1,n\}$. Thus
\[
 F(A)=(1,n,n,n),\quad F(B)=(1,n-1,n+1,n),\quad
 F(C)=(1,n,n-1,n+1).
\]
Equality in the sum of squares also implies that $v$ vanishes on every
remaining column. Such a column is therefore
\[
 (0,b-t,b,b+t),\qquad b\geq|t|,\quad b>0,
\]
and its central dimension is $3b\delta+2t$. Subtracting the four low columns
from $Q$ and comparing its entries gives
\begin{equation}\label{eq:quadratic-family-gram}
 \sum b^2=6n^2+2,\qquad \sum bt=2n,\qquad\sum t^2=2.
\end{equation}
Thus precisely two high columns have nonzero $t$, which equals $1$ or $-1$.
Set $T=3n\delta+2$. Subtracting the first twist-trace equation for $X$ from
that for $Z$, then dividing by two, gives
\begin{equation}\label{eq:quadratic-family-ttrace}
 \sum t(3b\delta+2t)\theta=-T.
\end{equation}
If the two signs of $t$ differ, their $b$ values differ by $2n$, by
\eqref{eq:quadratic-family-gram}. Their positive dimensions then differ by
$6n\delta+4=2T$, contradicting the reverse triangle inequality in
\eqref{eq:quadratic-family-ttrace}. Both signs cannot be negative. Hence both
are positive; write their $b$ values as $p,q$, with $p+q=2n$.
Their dimensions $\alpha=3p\delta+2$, $\beta=3q\delta+2$ satisfy
\[
 \alpha+\beta=2T,\qquad
 \alpha\beta=(9pq+4)\delta^2,\qquad T^2=(9n^2+4)\delta^2.
\]
If their twists are $\zeta,\eta$, equation
\eqref{eq:quadratic-family-ttrace} implies
\[
 \zeta\eta^{-1}+\eta\zeta^{-1}
 =2-\frac{3(9n^2+4)}{9pq+4}.
\]
This rational number is an algebraic integer and lies in $[-2,2]$.
Since $pq\leq n^2$, it is at most $-1$. The value $-2$ would require
$36pq=27n^2-4$, which is impossible modulo nine. The value $-1$ forces
$pq=n^2$, hence $p=q=n$. The two twists then sum to $-1$ and are the
primitive cube roots of unity.

All other columns have $t=0$ and are of the form $(0,b,b,b)$. Their squared
multiplicities sum to $4n^2+2$. Write $S=\sum b^2\theta$ over these uniform
columns; it is a sum of $4n^2+2$ roots of unity, with repetitions.
The first twist-trace equation for $X$ gives
\[
 S=n^2-n\delta.
\]
For completeness, the low contribution is
$n(1+3n\delta)+(2n-1)T$; the two exceptional high columns contribute
$-(n-1)T$, so the equation is
$0=n(1+3n\delta)+nT+3\delta S$, which is equivalent to the claimed identity
using $\delta^{-1}=\delta-3n$.

Writing $9n^2+4=t^2d$ with $d$ squarefree, we have
\[
 S=-\frac{n^2}{2}-\frac{nt}{2}\sqrt d.
\]
Lemma~\ref{lem:odd-quadratic-length} therefore requires
\[
 4n^2+2\geq\frac{n^2+nt\lambda}{2}.
\]
Condition~\eqref{eq:rootlength-parameter} contradicts this inequality.
For the finite parameters $3\leq n\leq127$, the supplement verifies it
exactly from the squarefree part $d$ of $9n^2+4$, its conductor and
$\lambda$. These are scalar integer and rational calculations, not a
search for roots of unity. The uniform estimate valid for every $n>2$
is proved in the companion paper; it is not a premise of the finite
certificate.

For $n=2$, the same reduction gives eighteen roots whose sum is
$S_1=-2-2\sqrt{10}$. The second indicator of $X$ is
$\epsilon\in\{1,-1\}$. The same twist-trace equation with twists squared,
using that the two exceptional twists still sum to $-1$, gives
\[
 S_2=S_1+\epsilon\frac{D}{3\delta}
     =-2+(2\epsilon-2)\sqrt{10}.
\]
For $\epsilon=-1$, Lemma~\ref{lem:rootlength40} applied to the eighteen
squared roots requires at least $34$ roots. For $\epsilon=1$, the same
lemma applied to $S_1$ shows that the normalized absolute trace of $S_2$
is $2$, whereas $S_2=-2$. Both cases are impossible.

Finally, a spherical categorification has a dimension character with no zero entries,
so the displayed character table restricts it to the two quadratic characters.
These are Galois conjugate; conjugating its structure constants if necessary
would give a pseudo-unitary categorification. This completes the proof.
\end{proof}
